\documentclass[11pt]{article}
\usepackage{mathblatt}
% Präambel des Setzers (werkzeuge/setzer-praeambel.tex, wortgleich); dazu \gteil (Muster T2U)
\makeatletter
\setlength{\mbpfbe}{0mm}% keine Punkte (pruefung.md § 5: nur beim Prüfstein)
\newlength{\gEin}\setlength{\gEin}{0pt}
\let\gSchrift\relax
\renewcommand{\pfkreuzzeile}[1]{\par\noindent\hangindent3.2em\hangafter1\hspace*{1.6em}$\square$\ #1\par}
\newcommand{\gkreuz}[1]{$\square$\ #1\hspace{2em}}
\newcommand{\gkreuzzeile}[2]{\par\noindent\hangindent2.2em\hangafter1\mbox{#1\,$\square$}\ #2\par}% Buchstabe, Option: eine je Zeile (neben Skizze)
% links, was man ansieht, rechts, was man tut; Spaltengrenze überall an derselben Stelle (gemeinsam.md „Satz“)
% v0.17: Fußeintrag „Nr: Wert“ – kurz in einer Box (nie getrennt), lang umbrechbar mit der Nummer am Anfang
% (vorher \mbox: lange Lösungen liefen über den Rand, AS40 138 pt); robust, weil \fusshilfe in die .aux schreibt
\NewDocumentCommand{\gfuss}{m m}{\sbox0{#1:~#2}\ifdim\wd0<0.45\textwidth\mbox{#1:~#2}\else\mbox{#1:}\nobreak\ #2\fi}
\newsavebox{\gbox}\newlength{\gLinks}
% v0.10 (Lehrer 13.10.): keine starre Halbierung – die linke Spalte ist so breit wie die Figur (+ 6 mm), die
% Antwort steht direkt daneben; ist die Figur breiter als 0,62 Zeile, steht die Antwort darunter
\newcommand{\gzwei}[2]{\par\smallskip\noindent\sbox{\gbox}{#1}%
  \setlength{\gLinks}{\dimexpr\wd\gbox+6mm\relax}%
  \ifdim\gLinks>0.62\linewidth
    \usebox{\gbox}\par\smallskip\noindent\begin{minipage}{\linewidth}#2\end{minipage}%
  \else
    \begin{minipage}[t]{\gLinks}\vspace{0pt}\usebox{\gbox}\end{minipage}%
    \begin{minipage}[t]{\dimexpr\linewidth-\gLinks\relax}\vspace{0pt}\raggedright #2\end{minipage}%
  \fi\par}
% v0.9 (12.10., pruefung.md § 5 „Marke mit B, Bedingungen rechts“): drittes Argument von pfaufg = Bedingungen
% der Prüfung (v0.10: \pfohnetr = TR durchgestrichen, „FOR“), klein grau im rechten Rand auf der Grundlinie der Nummer, Zeilen mit \\
\newcommand{\gbedrand}[2]{\ifblank{#2}{}{\rlap{\hspace*{\dimexpr#1+2mm\relax}%
  \parbox[t]{15mm}{\raggedright\scriptsize\leavevmode\color{mbgrau}\strut#2}}}}
% Teil eines Bündels (a, b …): Marke links im Rand, Bedingungen rechts im Rand, auf der Zeile des Teils
\newcommand{\gteilrand}[2]{\llap{\mbpfrandmarke{#1}\hspace*{\mbpfnr}}\gbedrand{\linewidth}{#2}}
% Teil eines Bündels: Buchstabe in eigener Spalte, Inhalt eingerückt (\gEin = Einzug, damit die Spaltengrenze von
% \gzwei an derselben Stelle bleibt)
\newenvironment{gbteil}[3]{\par\smallskip\noindent\gteilrand{#2}{#3}\makebox[\mbpfteil][l]{\textbf{#1}}%
  \begin{minipage}[t]{\dimexpr\linewidth-\mbpfteil\relax}\setlength{\gEin}{\mbpfteil}\raggedright\strut\ignorespaces}%
 {\end{minipage}\par\addvspace{6pt}}% v0.18 (PG80 Nr. 1): Abstand, damit die Linie von a) nicht an b) rutscht
% Grafik höchstens so breit wie die Spalte (Bündel nebeneinander, Koordinatensystem ohne Karo)
\newsavebox{\gpbox}
\newcommand{\gpassend}[1]{\sbox{\gpbox}{#1}\ifdim\wd\gpbox>\linewidth\typeout{GPASSEND: verkleinert \the\wd\gpbox\space auf \the\linewidth}\resizebox{\linewidth}{!}{\usebox{\gpbox}}\else\usebox{\gpbox}\fi}
% Skizzenfeld (v0.10, Lehrer 13.10.): Karo ab dem linken Rand über die Zeile, „Skizze“ klein grau oben links im Feld
% (wie \pfskizzenfeld der Vorlage); der rechte Rand bleibt den Bedingungen
\newcommand{\gskizze}[1]{\par\smallskip\noindent\begin{tikzpicture}
  \clip (0,0) rectangle ({\linewidth-1mm},{#1*10mm});
  \draw[black!16,line width=0.3pt,step=5mm] (0,0) grid ({\linewidth},{#1*10mm});
  \node[anchor=north west,inner sep=2pt,font=\scriptsize,text=mbgrau,fill=white] at (0,{#1*10mm}) {Skizze};
  \end{tikzpicture}\par}
\RenewDocumentEnvironment{pfaufg}{m m m}{%
  \begin{lrbox}{\mbaufgabebox}\begin{minipage}[t]{\linewidth}%
  \gSchrift\noindent\hspace*{\gEin}\mbpfrandmarke{#1}\gbedrand{\linewidth-\gEin}{#3}%
  \begin{minipage}[t]{\mbpfnr}\strut\textbf{#2}\end{minipage}%
  \begin{minipage}[t]{\dimexpr\linewidth-\gEin-\mbpfnr\relax}\raggedright\strut\ignorespaces}%
 {\end{minipage}%
  \end{minipage}\end{lrbox}%
  \setlength{\mbaufgabehoehe}{\dimexpr\ht\mbaufgabebox+\dp\mbaufgabebox\relax}%
  \ifdefempty{\gkopfnach}{\par\addvspace{5pt}\Needspace*{\mbaufgabehoehe}}{%
    % Stufenkopf, der mit dieser (ersten) Aufgabe zusammenbleibt: Platz für beide verlangen, dann den Kopf setzen
    % (12.10., Quadratische: ein Kopf stand allein unten, die Aufgabe mit Koordinatensystem auf der nächsten Seite)
    % v0.17 (AS40 „Terme deuten“ allein unten): hinter dem Kopf steht das Sprungziel (\special, ein Knoten, nach
    % dem die Seite brechen durfte) – darum \nobreak nach dem Kopf; Vorhalt 56 pt = Kopf mit Linie und Abständen
    \par\Needspace*{\dimexpr\mbaufgabehoehe+56pt\relax}\gkopfnach\par\nobreak\gdef\gkopfnach{}\addvspace{5pt}}%
  \noindent\usebox{\mbaufgabebox}\par\addvspace{5pt}}
% v0.18 (plan.md § 10, 15.10. l): Abschnittsbuchstabe A, B, C vor dem Namen im Kopf – als Vorsatz, damit der
% Name im TeX gleich bleibt (Bestellkarte, Sprungziele lesen ihn)
\def\gbstakt{}\newcommand{\gbstvor}[1]{\gdef\gbstakt{#1}}
\NewCommandCopy\gpfstufekopfalt\pfstufekopf
\RenewDocumentCommand{\pfstufekopf}{o m m}{\IfValueTF{#1}{\gpfstufekopfalt[#1]}{\gpfstufekopfalt}%
  {\ifx\gbstakt\empty\else\gbstakt\enspace\fi #2}{#3}\gdef\gbstakt{}}
% v0.18 (16.10. e): Begriff grau, ohne Nummer, ohne Rahmen über der ersten Aufgabe der Stufe
\newcommand{\gbegriff}[1]{\par\nobreak\noindent\hspace*{\mbpfnr}\parbox[t]{\dimexpr\linewidth-\mbpfnr\relax}{\raggedright\small\color{mbgrau}#1}\par\nobreak\addvspace{3pt}}
% v0.18 (16.10. g): Lösungsblatt – Kopf mit kräftigem Strich, Abschnitt klein grau als Zwischenzeile
\newcommand{\glkopf}[1]{\noindent{\large\bfseries Lösungen \textperiodcentered{} #1}\par\vspace{3pt}\hrule height 1.2pt\par\medskip}
\newcommand{\glzwischen}[1]{\par\addvspace{8pt}\Needspace*{4\baselineskip}\noindent{\footnotesize\color{mbgrau}#1}\par\nobreak\vspace{1pt}}
% v0.18 (15.10. o, w): Blattende – Kern ohne Rahmen, dann graue Zeile „Mehr: n Aufgaben“
\newcommand{\gkern}[1]{\par\addvspace{14pt}\Needspace*{6\baselineskip}\noindent{\bfseries Kern}\par\nobreak\smallskip
  {\small\raggedright #1\par}}
\newcommand{\gmehr}[1]{\par\addvspace{10pt}\noindent{\footnotesize\color{mbgrau}Mehr: #1 Aufgaben}\par}
% Stufenkopf vormerken: \gkopfvor{<Kopf>} – die nächste pfaufg setzt ihn, sobald beide zusammen Platz haben
\def\gkopfnach{}
\newcommand{\gkopfvor}[1]{\gdef\gkopfnach{#1}}
% Teilaufgabe mit Skizze links in gzwei: Buchstabe in eigener Spalte, darunter Text und Skizze
\newcommand{\gteil}[2]{\makebox[\mbpfteil][l]{\textbf{#1}}\parbox[t]{\dimexpr0.5\textwidth-\gEin-\mbpfnr-5mm-\mbpfteil\relax}{\raggedright #2}}
% Stufenkopf der Vorlage (mathblatt.sty 2026-10-12: Linie, Abstand, 13 pt, \pfjahre) ohne ihren Vorhalt von
% 14 Zeilen: der Setzer fordert vor dem Kopf 0,16 Seitenhöhe (Kopf + erste Aufgabe), mehr nicht
\setlength{\mbpfstufevorhalt}{0pt}
\makeatother
\newcommand{\karo}[1]{\par\smallskip\noindent\begin{tikzpicture}
  \clip (0,0) rectangle ({\linewidth-1mm},{#1*10mm});
  \draw[black!16,line width=0.3pt,step=5mm] (0,0) grid ({\linewidth},{#1*10mm});
  \end{tikzpicture}\par}
\newcommand{\kk}[1]{$\square$\,#1\hspace{0.9em}}
\newcommand{\linie}[1]{\,{\color{black!45}\rule[-1pt]{#1}{0.4pt}}\,}
\newcommand{\pkt}{\textperiodcentered{}}
% v0.14 (14.10.): Stufenkopf mit fünf Punkten (Option = Jahre für die Bestellkarte) oder „zuletzt 19“
\NewDocumentCommand{\gpunkte}{o m}{#2}
\newcommand{\gzuletzt}[1]{zuletzt #1}
% v0.7 (12.10.): Schülerfrage unter dem Stufenkopf; Ankreuzoption als Bild (Buchstabe und Kästchen darunter)
\newcommand{\gfrage}[1]{\par\nobreak\noindent\hspace*{\mbpfnr}{\small\itshape #1}\par\nobreak\addvspace{2pt}}
\newlength{\gbildbreite}\setlength{\gbildbreite}{0.32\linewidth}% Setzer setzt sie je Reihe (vier Bilder: 0,24)
\newcommand{\gbildoption}[3]{\begin{minipage}[t][#3][b]{\gbildbreite}\centering #2\par\smallskip\mbox{#1\,$\square$}\end{minipage}}
% Lineare (12.10.): Bild mit Buchstabe ohne Kästchen (Diagramm zuordnen); Aussage mit „wahr / falsch“ zum Ankreuzen
\newcommand{\gbildname}[3]{\begin{minipage}[t][#3][b]{\gbildbreite}\centering #2\par\smallskip\mbox{\textbf{#1}}\end{minipage}}
\newcommand{\gwahrfalsch}[2]{\par\noindent\hangindent2.2em\hangafter1\makebox[2.2em][l]{(#1)}#2\hfill\mbox{wahr $\square$\quad falsch $\square$}\par}
\begin{document}

\pfheftstil
\pflaufkopf[Abitur GK \pkt{} Prüfung]{Dreieck aus Tangente und Normale}{AA90}\pfseitevon
\pfheftkopf{Dreieck aus Tangente und Normale}{}
\pfstufenzeile{Umfang oder Fläche gesucht \pkt{} Normale begründen und deuten}

% ===================== Stufe: Umfang oder Fläche gesucht =====================
\gkopfvor{\pfstufekopf{Umfang oder Fläche gesucht}{\gpunkte[26 24 23 22]{\textcolor{black!80}{$\bullet$}\,\textcolor{black!25}{$\circ$}\,\textcolor{black!80}{$\bullet$}\,\textcolor{black!80}{$\bullet$}\,\textcolor{black!80}{$\bullet$}}}\special{pdf:dest (stufe-umfang-oder-flache-gesucht) [@thispage /XYZ @xpos @ypos null]}}
% kein Einstieg: keine Schülerfrage in der Gliederung (Schülerfrage, -Fälle, -Form, -Weg, -Skizze; oder Frage) – Einstieg nur von dort, nie aus der Bank (pruefung.md § 5, 12.10.)
% 2023-bebb-gk-A1.1b-m1 (Original 2023-bebb-gk-A1.1b, Schwere II, rang 32; herausgelöste Fassung: keine)
\begin{pfaufg}{GK~’23\\A1.1b}{1.}{\pfohnetr}
Die Tangenten $t_1$: $y = \tfrac{3}{4}x + 3$ und $t_2$: $y = -\tfrac{3}{4}x + 3$ schließen mit der $x$-Achse ein Dreieck ein. Bestimme seinen Umfang.
\karo{4}
\fusshilfe{\gfuss{1}{18}}
\end{pfaufg}
% 2026-bb-gk-B2.2c-m1 (Original 2026-bb-gk-B2.2c, Schwere II, rang 33; herausgelöste Fassung: keine)
\begin{pfaufg}{GK~’26\\B2.2c}{2.}{}
Die Tangente $t$: $y = 2x + 5$ bildet mit den Koordinatenachsen ein Dreieck. Berechne seinen Umfang.
\karo{4}
\fusshilfe{\gfuss{2}{13,09}}
\end{pfaufg}
% 2022-bebb-gk-B2.1g-m1 (Original 2022-bebb-gk-B2.1g, Schwere II, rang 34; herausgelöste Fassung: keine)
\begin{pfaufg}{GK~’22\\B2.1g}{3.}{}
Gegeben ist die in $\mathbb{R}$ definierte Funktion $f$ mit $f(x) = (x + 3) \cdot e^{-x}$ und $f'(x) = -(x + 2) \cdot e^{-x}$. Die Tangente an den Graphen von $f$ im Punkt $(0 \mid f(0))$ bildet mit den Koordinatenachsen ein Dreieck. Berechne den Umfang dieses Dreiecks. Gib die Koordinaten des Punkts an, der von allen drei Eckpunkten gleich weit entfernt ist.
\karo{4}
\fusshilfe{\gfuss{3}{7,85; Mitte der Hypotenuse (0,75|1,5)}}
\end{pfaufg}
% 2024-bebb-gk-B2.2d-m1 (Original 2024-bebb-gk-B2.2d, Schwere II, rang 35; herausgelöste Fassung: keine)
\begin{pfaufg}{GK~’24\\B2.2d}{4.}{}
Gegeben ist die Funktion $f$ mit $f(x) = x^4 - 6x^2 + 5$, $x \in \mathbb{R}$; ihr Graph ist $G_f$. Die Gerade $g$ verläuft durch $A(-1 \mid f(-1))$ und steht senkrecht auf der Tangente $t$ an $G_f$ in $A$. Die Geraden $g$, $t$ und die $y$-Achse schließen ein Dreieck ein. Bestimme seinen Flächeninhalt.
\karo{4}
\fusshilfe{\gfuss{4}{$-\tfrac{1}{8}$ x $-$ $\tfrac{1}{8}$; 4,06}}
\end{pfaufg}

% ===================== Stufe: Normale begründen und deuten =====================
\gkopfvor{\pfstufekopf{Normale begründen und deuten}{\gpunkte[24 23 22]{\textcolor{black!25}{$\circ$}\,\textcolor{black!25}{$\circ$}\,\textcolor{black!80}{$\bullet$}\,\textcolor{black!80}{$\bullet$}\,\textcolor{black!80}{$\bullet$}}}\special{pdf:dest (stufe-normale-begrunden-und-deuten) [@thispage /XYZ @xpos @ypos null]}}
% kein Einstieg: keine Schülerfrage in der Gliederung (Schülerfrage, -Fälle, -Form, -Weg, -Skizze; oder Frage) – Einstieg nur von dort, nie aus der Bank (pruefung.md § 5, 12.10.)
% Teile eines Originals (2023-bebb-gk-B2.1): 2023-bebb-gk-B2.1g-m1, 2023-bebb-gk-B2.1h-m1
\begin{pfaufg}{}{5.}{}
% 2023-bebb-gk-B2.1g-m1 (Original 2023-bebb-gk-B2.1g, Schwere III, rang 36; herausgelöste Fassung: keine)
\vspace*{-\baselineskip}\vspace*{-\smallskipamount}
\begin{gbteil}{a)}{\raisebox{\dimexpr-\height+\ht\strutbox\relax}{\shortstack[r]{GK~’23\\B2.1g}}}{}
Gegeben ist die in $\mathbb{R}$ definierte Funktion $f$ mit $f(x) = \tfrac{1}{3}(x^2 - 9) \cdot e^{\tfrac{x}{2}}$ und $f'(x) = \left(\tfrac{1}{6}x^2 + \tfrac{2}{3}x - \tfrac{3}{2}\right) \cdot e^{\tfrac{x}{2}}$. Die Abbildung zeigt einen Teil des Graphen $G$ mit dem Wendepunkt $W(0{,}12 \mid -3{,}19)$; dort ist $f'(0{,}12) \approx -1{,}51$. Die Tangente an $G$ in $W$ begrenzt mit den Koordinatenachsen ein Dreieck; die zu ihr senkrechte Gerade durch $W$ begrenzt mit den Achsen ein weiteres Dreieck. Zeichne beide Dreiecke in die Abbildung ein. Begründe mithilfe der Zeichnung, dass das eine Dreieck eine maßstäbliche Vergrößerung des anderen ist.
\par\smallskip\noindent\begin{abiksys}[xmin=-3,xmax=6,ymin=-4,ymax=1,le=1]\funktionab[-0.4]{(\x^2-9)*exp(\x/2)/3}{G}{-0.7}{0.9}\wendepunkt{0.123}{-3.185}{W}\end{abiksys}\par
\karo{2}
\end{gbteil}
% 2023-bebb-gk-B2.1h-m1 (Original 2023-bebb-gk-B2.1h, Schwere III, rang 37; herausgelöste Fassung: keine)
\begin{gbteil}{b)}{\raisebox{\dimexpr-\height+\ht\strutbox\relax}{\shortstack[r]{GK~’23\\B2.1h}}}{}
Gegeben ist die in $\mathbb{R}$ definierte Funktion $f$ mit $f(x) = \tfrac{1}{3}(x^2 - 9) \cdot e^{\tfrac{x}{2}}$; ihr Graph ist $G$. Es gibt eine Zahl $a$ mit $f'(a) = -\tfrac{1}{f''(a)}$. Erläutere die geometrische Bedeutung dieser Gleichung.
\karo{3}
\end{gbteil}
\fusshilfe{\gfuss{5}{a)~Tangente ($-$1,99|0) und (0|$-$3,00), Normale (4,92|0) und (0|$-$3,27); ähnlich (Streckfaktor $\approx$ 1,6); b)~die Tangente an G in a steht senkrecht auf der Tangente an den Graphen von f' in a}}
\end{pfaufg}
% 2024-bebb-gk-B2.1d-m1 (Original 2024-bebb-gk-B2.1d, Schwere III, rang 38; herausgelöste Fassung: keine)
\begin{pfaufg}{GK~’24\\B2.1d}{6.}{}
Gegeben ist die in $\mathbb{R}$ definierte Funktion $f$ mit $f(x) = (x - 3) \cdot e^{-\tfrac{x}{3} + 2}$ und $f'(x) = \left(2 - \tfrac{x}{3}\right) \cdot e^{-\tfrac{x}{3} + 2}$. Die Abbildung zeigt $G_f$ und $G_{f'}$ ohne Skalierung. Der Graph von $f$ hat den Wendepunkt $(9 \mid 6e^{-1})$. Entscheide, ob die Aussagen wahr sind, und begründe: (I) Die Gerade durch $A(3 \mid 0)$ und $B(9 \mid f(9))$ steht senkrecht auf der Tangente mit dem stärksten Gefälle. (II) Jede Parallele zur $x$-Achse mit der Gleichung $y = k$, $k \le -e^{-1}$, hat mit $G_{f'}$ keinen gemeinsamen Punkt.
\gzwei{\begin{abiksys}[xmin=1,xmax=16,ymin=-2,ymax=4.5,le=0.4,yle=0.6,ohnezahlen,ohnegitter]\funktion[13]{(\x-3)*exp(-\x/3+2)}{G_f}\funktion[14]{(2-\x/3)*exp(-\x/3+2)}{G_{f'}}\end{abiksys}}{\karo{4}}
\fusshilfe{\gfuss{6}{falsch; falsch}}
\end{pfaufg}
% 2022-bebb-gk-B2.2j-m1 (Original 2022-bebb-gk-B2.2j, Schwere III, rang 39; herausgelöste Fassung: keine)
\begin{pfaufg}{GK~’22\\B2.2j}{7.}{}
Der Querschnitt eines Betonteils (1 LE = 1 m, Abbildung) ist oben begrenzt durch den Graphen von $h$ mit $h(x) = -\tfrac{1}{16}x^3 + \tfrac{3}{8}x^2$ von $A(0 \mid 0)$ bis $B(4 \mid 2)$ und unten durch die Gerade $g$ mit $g(x) = \tfrac{1}{2}x - 1$ von $C(0{,}4 \mid -0{,}8)$ bis $D(4 \mid 1)$. Geradlinige Bohrungen sollen den oberen und den unteren Rand jeweils im rechten Winkel schneiden. Erläutere, wie man die Lage dieser Bohrungen ermitteln kann.
\gzwei{\begin{abiksys}[xmin=-0.5,xmax=4.5,ymin=-1.5,ymax=2.5,le=1]\funktionab[2.6]{-\x^3/16+3*\x^2/8}{h}{0}{4}\draw[line width=0.8pt] (0,0)--(0.4,-0.8)--(4,1)--(4,2);\node[above left,font=\small] at (0,0) {$A$};\node[above,font=\small] at (4,2) {$B$};\node[below,font=\small] at (0.4,-0.8) {$C$};\node[right,font=\small] at (4,1) {$D$};\end{abiksys}}{\karo{4}}
\fusshilfe{\gfuss{7}{$\tfrac{1}{2}$ lösen (x0 = 2 ± $\tfrac{2}{3}$ $\sqrt{3}$ $\approx$ 0,85 und 3,15); Gerade mit Steigung $-$2 durch (x0|h(x0))}}
\end{pfaufg}

\pfblattende{10}{26 & 25 & 24 & 23 & 22 & 21 & 20 & 19 & 18 & 17}{B2.2c &  & B2.1d B2.2d & A1.1b B2.1g B2.1h & B2.1g B2.2j & B2.2f &  &  &  & }%
{\pfausschnitt{Tangente \pkt{} Winkel am Graphen}{Dreieck aus Tangente und Normale}{}}
\end{document}